% Propietario conservado: /Users/ruben/Documents/New project/output/REVISION_ARGUMENTAL_APERTURAS_CIERRES_20260915/VIII/spanish_source/gestion/lectura_generada/05_CONSTRUCCION_CORRELATIVA_DEL_CONTINUO.tex
\section{Construction corrélative du continu : histoires, transport, incidence et terminal}
\label{viii:lectura:manuscrito-05}
L'arithmétique publiée par une forme normale conserve une lecture précise de l'état qui la produit. Pour situer cette lecture dans HMT, il est nécessaire de maintenir aussi les histoires, les feuillets et les opérations qui agissent sur eux. Ce chapitre réunit ce substrat commun. Ses différentes constructions s'obtiennent à partir du même arbre d'exécutions APP–TRIT–TPK ; les découpages de l'exposé permettent d'étudier leurs propriétés sans les transformer en domaines indépendants choisis par une application ultérieure.

Le point de départ est le noyau formel exposé précédemment. Une émission admissible contient l'opération d'APP, l'orientation du TRIT, le changement de feuillet, le résidu et le quotient, ainsi que la mise à jour effective de la mémoire du TPK. Lorsque seul un résidu est publié, l'histoire de l'exécution appartient toujours à l'objet de départ. Deux exécutions de lecture finale identique ne sont pas identifiées pour ce seul motif.

\subsection{Histoires admissibles et mémoire par composition}
Soit \(X_k\) l’ensemble des histoires admissibles de longueur \(k\). Ses éléments sont construits par émissions successives depuis le domaine initial du noyau. Nous écrivons \(h\varepsilon\) pour le prolongement de \(h\) par une émission admissible \(\varepsilon\), et \(\rho_k(h\varepsilon)=h\) pour la suppression de la dernière émission. L’étiquette de l’émission est conservée, même lorsque des étiquettes différentes produisent un même état publié.

Dans les deux feuillets arithmétiques, les données locales satisfont les divisions exactes

\[
\delta_\diamond=r_\diamond+9q_\diamond,
\qquad 0\leq r_\diamond<9,
\qquad \diamond\in\{\Sigma,\Pi\}.
\]
La coordonnée orientée du TRIT admet de même une écriture \(\Delta=o_r+3\kappa\), avec le représentant orienté déclaré par le noyau. La conservation du quotient permet de composer ces publications sans remplacer une émission par son résidu isolé.

Dans une carte de mémoire, la mise à jour d'une émission s'écrit

\[
U_\varepsilon(m)=C_\varepsilon m+\kappa_\varepsilon.
\]
Pour une trajectoire \(\alpha=(\varepsilon_1,\ldots,\varepsilon_k)\), l'opération effective est la composition ordonnée de ces mises à jour. Sa partie linéaire et sa translation sont

\[
C_\alpha=C_{\varepsilon_k}\cdots C_{\varepsilon_1},
\qquad
\kappa_\alpha=
\sum_{j=1}^{k}C_{\varepsilon_k}\cdots
C_{\varepsilon_{j+1}}\kappa_{\varepsilon_j}.
\]
Le produit vide de la dernière somme est l'identité.

\HMTresultado{Proposition}{Compatibilité affine de la concaténation}{rz-num-05-1}
Si \(\alpha\) précède \(\beta\), alors

\[
C_{\beta\alpha}=C_\beta C_\alpha,
\qquad
\kappa_{\beta\alpha}=C_\beta\kappa_\alpha+\kappa_\beta,
\qquad
U_{\beta\alpha}=U_\beta U_\alpha.
\]
\textbf{Démonstration.} Substituer \(U_\alpha(m)\) dans \(U_\beta\) produit \(C_\beta C_\alpha m+C_\beta\kappa_\alpha+\kappa_\beta\). La décomposition en partie linéaire et translation donne les trois identités. L'associativité de la composition démontre que le résultat ne dépend pas d'une subdivision auxiliaire de la trajectoire. \hfill\(\square\)

Cette identité de cocycle est le contenu précis de la mémoire accumulée. La quantité transportée par une seconde trajectoire dépend de la mise à jour de carte qu'elle réalise ; en général, ce n'est pas la somme sans transport de deux registres.

Le retour nonadique conserve une seconde distinction. Dans les coordonnées de phase et de tour, avec \(\bar r\in\{0,\ldots,8\}\), sa publication est

\[
\sigma_9(w,r)=
\left(w+\left\lfloor\frac{\bar r+1}{9}\right\rfloor,
r+1\pmod9\right),
\qquad
\sigma_9^9(w,r)=(w+1,r).
\]
Pendant neuf étapes, il se produit exactement un franchissement de \(8\) à \(0\). La phase revient et la coordonnée de tour augmente. L'holonomie enrichie retient en outre les mises à jour de mémoire qui accompagnent ce parcours.

\subsection{Raffinement de l'arbre et mesure des histoires}
On fixe l’ensemble initial fini non vide de germes admissibles, avec une masse initiale strictement positive et de somme \(1\). Dans ce paragraphe, chaque histoire a un nombre fini et positif de prolongements. La suppression d’émissions définit les applications de raffinement. L’admissibilité du noyau fournit les prolongements ; lorsqu’une résidence utilise une émission neutre, celle-ci donne explicitement une section de \(\rho_k\). L’existence de cette section se rapporte à l’arbre qui comprend cette émission, non à toute restriction arbitraire de routes.

Soit \(d(h)\geq1\) le nombre d'émissions admissibles en \(h\). Le choix équiprobable parmi ses enfants définit

\[
\mu(h\varepsilon)=\frac{\mu(h)}{d(h)}.
\]
Plus généralement, une distribution locale positive \(p(\varepsilon\mid h)\), dont la somme est \(1\), définit \(\mu(h\varepsilon)=\mu(h)p(\varepsilon\mid h)\). Cette seconde formulation permet de conserver les pondérations du lecteur lorsque les nombres de prolongements ne sont pas uniformes.

\HMTresultado{Proposition}{Compatibilité cylindrique et espérance conditionnelle}{rz-num-05-2}
À chaque niveau, on a

\[
\sum_{\rho_k(y)=x}\mu(y)=\mu(x).
\]
Dans les espaces \(L^2(X_k,\mu_k)\), les applications

\[
(J_kf)(y)=f(\rho_k(y)),
\qquad
(E_kg)(x)=\sum_{\rho_k(y)=x}\frac{\mu(y)}{\mu(x)}g(y)
\]
satisfont \(E_k=J_k^*\), \(E_kJ_k=I\) et \(\|J_kf\|=\|f\|\).

\textbf{Démonstration.} La première identité est la normalisation des probabilités locales. En l'appliquant à \(|f(x)|^2\), on obtient l'isométrie. En regroupant par parents la somme qui définit \(\langle J_kf,g\rangle\), on obtient la formule de \(E_k\) ; sur une fonction constante sur chaque fibre, cette formule renvoie sa valeur. \hfill\(\square\)

Les masses compatibles définissent une mesure de probabilité sur l’espace des histoires infinies : le cylindre d’une histoire finie reçoit sa masse \(\mu(h)\). La mesure s’étend à partir de l’algèbre des cylindres, et son unicité découle du fait que ceux-ci engendrent la tribu. Si l’arbre est à ramifications finies et si chaque histoire possède un prolongement, la limite projective de ses niveaux par suppression d’émissions est non vide. Dans la topologie cylindrique, chaque cylindre non vide est de mesure positive.

Cette mesure distingue les histoires qui convergent vers une même publication. La mesure uniforme sur les valeurs terminales serait une lecture différente, car elle éliminerait les multiplicités que la construction vient de conserver.

\subsection{Transport par feuillets et comptabilité exacte des sorties}
Considérons d'abord l'espace dont la base orthonormée est formée des histoires du niveau \(k\). Le transport qui conserve l'étiquette de chaque prolongement est

\[
\mathcal U_ke_h=
\sum_{\varepsilon\ \mathrm{admisible\ en}\ h}
\sqrt{p(\varepsilon\mid h)}\,e_{h\varepsilon}.
\]
Les ensembles d’enfants de parents distincts sont disjoints. La normalisation de \(p\) démontre donc \(\mathcal U_k^*\mathcal U_k=I\). La représentation au moyen des masses d’histoires est reliée à celle-ci par le changement de base diagonal qui incorpore \(\sqrt{\mu(h)}\).

Soient \(P_k,Q_k\) les projecteurs complémentaires sur les deux feuillets. Nous séparons la conservation de feuillet et l'échange de feuillet :

\[
A_k=P_{k+1}\mathcal U_kP_k+Q_{k+1}\mathcal U_kQ_k,
\qquad
\eta_k=Q_{k+1}\mathcal U_kP_k+P_{k+1}\mathcal U_kQ_k.
\]
\HMTresultado{Proposition}{Identité de Gram par feuillets}{rz-num-05-3}
Pour un transport linéaire \(U\), isométrique ou non, la décomposition précédente satisfait

\[
A^*A+\eta^*\eta=PU^*UP+QU^*UQ.
\]
En particulier, pour \(U=\mathcal U_k\),

\[
A_k^*A_k+\eta_k^*\eta_k=I.
\]
\textbf{Démonstration.} Dans \(A^*A\) et \(\eta^*\eta\), les termes mixtes qui contiennent \(P_{k+1}Q_{k+1}\) s’annulent. En additionnant les termes restants, on fait apparaître \(P_kU^*(P_{k+1}+Q_{k+1})UP_k\), et son analogue avec \(Q_k\). La première identité résulte de la complémentarité des projecteurs ; la seconde utilise en outre \(U^*U=I\). \hfill\(\square\)

La première identité est la formule générale. Remplacer son membre de droite par \(U^*U\) requiert l'annulation des blocs croisés de ce gramien. La conservation isométrique du transport des histoires fournit ici cette condition ; elle ne doit pas être attribuée automatiquement à une compression ultérieure.

Écrivons \(F_{0,0}=I\), \(F_{0,n}=A_{n-1}\cdots A_0\) et \(T_N=F_{0,N}^*F_{0,N}\).

\HMTresultado{Théorème}{Décomposition finie avec résidu terminal}{rz-num-05-4}
Pour chaque horizon \(N\),

\[
I=\sum_{n=0}^{N-1}
F_{0,n}^*\eta_n^*\eta_nF_{0,n}+T_N.
\]
La suite \(T_N\) est positive et décroissante, et possède une limite forte positive \(T_\infty\). La somme des sorties converge fortement vers \(I-T_\infty\).

\textbf{Démonstration.} La proposition~\ref{viii:rz-num-05-3} donne \(T_n-T_{n+1}=F_{0,n}^*\eta_n^*\eta_nF_{0,n}\). La somme télescopique prouve l’égalité finie et la monotonie. Pour chaque vecteur, les formes quadratiques décroissent vers une limite ; la polarisation définit une forme bornée positive et, par le théorème de représentation des formes bornées, un opérateur \(T_\infty\). La convergence est forte car, pour un opérateur positif \(S\leq I\), \(\|Sv\|^2\leq\langle v,Sv\rangle\) ; on l’applique à \(S=T_N-T_\infty\). \hfill\(\square\)

Le résidu terminal est conservé jusqu'à la démonstration de son extinction dans le domaine considéré. Cette distinction empêche qu'un retour de phase, à lui seul, se transforme en une hypothèse d'épuisement de toutes les histoires.

\subsection{Bilan orienté et annulation dans le raffinement}
Une histoire porte son orientation cumulée \(o(h)\) et les décomptes de ses visites aux deux feuillets. Le bilan local correspondant est

\[
b(h)=o(h)\bigl(N_\Sigma(h)-N_\Pi(h)\bigr),
\qquad
b^+=\frac{|b|+b}{2},\quad b^-=\frac{|b|-b}{2}.
\]
Lorsqu'un lecteur publie la moyenne sur les prolongements, on définit \(b_c=Eb_f\). L'orientation reste dans l'argument de \(b_f\) ; elle n'est pas éliminée avant le calcul de la moyenne.

\HMTresultado{Proposition}{Défaut exact d'annulation}{rz-num-05-5}
La quantité

\[
\kappa_{\mathrm{can}}=
\frac{E|b_f|-|Eb_f|}{2}
\]
est non négative et satisfait

\[
E(b_f^+)=(b_c)^++\kappa_{\mathrm{can}},
\qquad
E(b_f^-)=(b_c)^-+\kappa_{\mathrm{can}}.
\]
\textbf{Démonstration.} L'inégalité triangulaire pondérée donne \(|Eb_f|\leq E|b_f|\). En écrivant les parties positive et négative au moyen de \((|b|\pm b)/2\), on obtient les deux égalités. \hfill\(\square\)

L’information omise lorsque seul le bilan est publié est l’annulation commune des contributions opposées. Cette quantité est déterminée à partir des mêmes prolongements que ceux utilisés par le transport précédent. Le facteur \(1/9\) n’apparaît que lorsque la fibre contient neuf enfants de poids égaux ; l’espérance conditionnelle est la règle qui demeure valide pour une ramification et des pondérations générales.

\subsection{Incidence ternaire des feuillets et complexe rectangulaire}
Les deux feuillets déterminent le module \(V=\mathbb F_3e_\Sigma\oplus\mathbb F_3e_\Pi\). Ses quatre directions projectives sont les droites engendrées par \(e_\Sigma,e_\Pi,e_\Sigma+e_\Pi,e_\Sigma-e_\Pi\). Elles donnent six paires non ordonnées. Les huit vecteurs non nuls sont leurs représentants orientés. Ainsi, les dénombrements \(4,6,8\) appartiennent à des opérations distinctes sur un même module de feuillets.

Sur les six positions de paires, nous définissons

\[
(Ba)_{\{L,L'\}}=a_L+a_{L'},
\qquad
s(q)=\sum_{\{L,L'\}}q_{\{L,L'\}}.
\]
Chaque direction intervient dans trois paires ; d'où \(sB=0\) dans \(\mathbb F_3\). Le lecteur \(W(q)=\mathbf1_{s(q)=0}-1/3\) est invariant sous \(q\mapsto q+Ba\). Sur l'espace ambiant uniforme \(\mathbb F_3^6\), ses moments sont

\[
\mathbb EW=0,\qquad
\mathbb EW^2=\frac29,\qquad
\mathbb EW^3=\frac2{27}.
\]
En effet, \(s\) est surjective et chaque fibre contient \(3^5\) éléments. Le lecteur vaut \(2/3\) avec la probabilité \(1/3\) et \(-1/3\) avec la probabilité \(2/3\). Ces moyennes décrivent l'espace ambiant uniforme ; la distribution induite par un sous-ensemble d'histoires doit être calculée avec ses propres multiplicités, déjà conservées dans \(X_k\).

Pour les ports effectifs d'un niveau, soient \(I\) et \(J\) les ensembles non vides de ports de ses deux feuillets et \(A_N=\mathbb Z/3^N\mathbb Z\). Fixons des ports de référence \(i_0,j_0\). Les applications

\[
\kappa(u,v)_{ij}=u_i-v_j,
\qquad
(\delta w)_{ij}=w_{ij}-w_{ij_0}-w_{i_0j}+w_{i_0j_0}
\]
définissent la suite suivante, valable sur l'anneau fini \(A_N\) :

\[
0\longrightarrow A_N\longrightarrow
A_N^I\oplus A_N^J\mathrel{\mathop{\longrightarrow}^{\kappa}}
A_N^{I\times J}\mathrel{\mathop{\longrightarrow}^{\delta}}
A_N^{(I\setminus\{i_0\})\times(J\setminus\{j_0\})}
\longrightarrow0.
\]
\HMTresultado{Proposition}{Exactitude de l'incidence rectangulaire}{rz-num-05-6}
La suite précédente est exacte, sa première flèche étant donnée par la diagonale constante.

\textbf{Démonstration.} L’égalité \(u_i-v_j=0\) pour chaque paire impose que toutes les coordonnées de \(u\) et \(v\) soient une même constante. La substitution directe donne \(\delta\kappa=0\). Si \(\delta w=0\), prenons \(u_i=w_{ij_0}\) et \(v_j=w_{i_0j_0}-w_{i_0j}\). Alors \(u_i-v_j=w_{ij}\). Enfin, toute matrice du dernier terme se prolonge avec une ligne et une colonne de référence nulles ; son image par \(\delta\) est la matrice initiale. La preuve n’utilise que des sommes et des différences, et n’exige donc pas de division par une unité de l’anneau. \hfill\(\square\)

La réduction de \(A_{N+1}\) à \(A_N\) commute avec les deux applications. Les formules de reconstruction sont les mêmes à chaque profondeur. L’incidence et son exactitude sont ainsi compatibles avec le raffinement arithmétique de l’état de départ.

\subsection{Complexe de mémoire et compatibilité des incidences}
La même coordonnée accumulée de mémoire \(c_y\in\mathbb Z\), dans un lieu de résidence \(y\), définit un complexe libre entier. Sa réduction à \(A_N\) utilise le même registre entier avant de le réduire ; les profondeurs successives sont donc compatibles. Ses générateurs sont \(f_y\) en degré \(2\), \(e_y,b_y,r_y\) en degré \(1\), et \(g_y\) en degré \(0\) :

\[
\partial f_y=3c_y e_y+b_y,\qquad
\partial e_y=g_y,\qquad
\partial b_y=-3c_y g_y,\qquad
\partial r_y=0.
\]
L'identité \(\partial^2=0\) se vérifie sur \(f_y\) au moyen de \(3c_yg_y-3c_yg_y=0\), et elle est immédiate sur les autres générateurs.

Si une trajectoire \(\alpha:y\to y'\) met à jour la mémoire, son transport dans le complexe envoie les générateurs homonymes sur leurs nouveaux représentants, sauf

\[
\Psi_\alpha(b_y)=b_{y'}+3(c_{y'}-c_y)e_{y'}.
\]
\HMTresultado{Proposition}{Naturalité du complexe de mémoire}{rz-num-05-7}
On a \(\partial\Psi_\alpha=\Psi_\alpha\partial\) et \(\Psi_{\beta\alpha}=\Psi_\beta\Psi_\alpha\). Les applications se prolongent linéairement sur l'anneau de coefficients choisi.

\textbf{Démonstration.} Sur \(f_y\), l'image de son bord est \(3c_ye_{y'}+b_{y'}+3(c_{y'}-c_y)e_{y'} =3c_{y'}e_{y'}+b_{y'}\). Sur \(b_y\), le bord de son image est \(-3c_{y'}g_{y'}+3(c_{y'}-c_y)g_{y'}=-3c_yg_{y'}\). Les deux autres cas sont immédiats. Dans une concaténation, les différences de mémoire s'additionnent télescopiquement, ce qui prouve la seconde identité. \hfill\(\square\)

La relation entre deux représentants d'une même résidence peut également s'écrire comme un bord explicite. Pour une coordonnée entière \(v\), réduite dans le même anneau le cas échéant, soient

\[
z^-=r_y,\qquad z^+=r_y+3c_yv e_y+v b_y,
\qquad h_*=-vf_y.
\]
Alors \(\partial h_*=z^--z^+\). Sous un transport qui met aussi à jour \(v\), la correction \(K_\alpha=(v_{y'}-v_y)f_{y'}\) satisfait \(K_{\beta\alpha}=K_\beta+\Psi_\beta K_\alpha\). Les symboles \(K_\alpha\) de cette homotopie locale ne désignent pas le sceau dodécaphasique \(K\) : ils appartiennent à un autre domaine et leur fonction est spécifiée par l'identité précédente. Leur fonction peut aussi être vérifiée dans l'égalité

\[
z^+_{y'}-\Psi_\alpha z^+_y=\partial K_\alpha.
\]
Les deux membres valent \((v_{y'}-v_y)(3c_{y'}e_{y'}+b_{y'})\).

\HMTresultado{Proposition}{Homologie et mémoire du complexe local}{rz-num-05-8}
Sur \(\mathbb Z\), ou après réduction à \(A_N\), on a

\[
H_0(\Gamma)=H_2(\Gamma)=0,
\qquad H_1(\Gamma)\cong A_N[r_y]
\]
dans la réalisation modulaire, et \(H_1(\Gamma)\cong\mathbb Z[r_y]\) dans la réalisation entière.

\textbf{Démonstration.} On définit l’homotopie de degré \(1\) par \(h(g_y)=e_y\), \(h(b_y)=f_y\), et zéro sur les autres générateurs. Soient \(p\) la projection sur le terme engendré par \(r_y\) et \(\iota\) son inclusion. La vérification sur chaque générateur donne

\[
\partial h+h\partial=I-\iota p.
\]
En particulier, sur \(b_y\), les termes \(3c_ye_y\) et \(-3c_ye_y\) s’annulent. Le complexe se rétracte par homotopie sur le module engendré par \(r_y\), concentré en degré \(1\), ce qui démontre l’assertion. \hfill\(\square\)

Le transport \(\Psi_\alpha\) est un cisaillement de déterminant \(1\). Ainsi, ce complexe conserve explicitement la manière dont \(c_y\) change au niveau des chaînes, mais son homologie locale ne distingue pas les valeurs de cette coordonnée. La mémoire complète demeure dans les étiquettes et dans les transports de l'état source ; elle ne se restitue pas à partir de \(H_1\) isolément. Une torsion numérique dépendant de la mémoire requiert l'assemblage et la lecture déterminantale spécifiés, outre ce complexe local.

\subsection{Assemblage avec les chemins terminaux}
Pour un horizon fini, les résidences et leurs prolongements forment une catégorie de chemins. À chaque résidence \(\nu\) est attribué le module libre \(W(\nu)=A_N[\operatorname{Path}(\nu,\mathrm{Term})]\), qui conserve chaque continuation terminale comme générateur. Une trajectoire initiale agit dans \(W\) par précomposition. Le complexe \(\Gamma(\nu)\) de la section précédente agit de manière covariante au moyen de \(\Psi\).

L'assemblage utilise les deux transports dans le même module bigradué :

\[
\mathcal B_{q,r}=
\bigoplus_{\nu_0\to\cdots\to\nu_q}
W(\nu_q)\otimes_{A_N}\Gamma_r(\nu_0).
\]
On utilise le complexe normalisé des chemins : les chaînes dégénérées contenant des identités sont omises. L'horizon fini borne alors le nombre de prolongements stricts et le degré horizontal. Le bord horizontal supprime un lieu de résidence intermédiaire en composant les flèches adjacentes. À la première extrémité, il transporte \(\Gamma\) ; à la dernière, il précompose \(W\). Avec les signes alternés, on obtient le bord des chemins \(d_{\mathrm{bar}}\). La différentielle totale est

\[
D=d_{\mathrm{bar}}+(-1)^q\partial_\Gamma.
\]
Les identités de faces annulent par paires les termes de \(d_{\mathrm{bar}}^2\). La naturalité de la proposition~\ref{viii:rz-num-05-7} implique que \(d_{\mathrm{bar}}\) commute avec \(\partial_\Gamma\) ; le facteur \((-1)^q\) annule les termes croisés du carré. Par conséquent, \(D^2=0\). La compatibilité se démontre sur les générateurs de chemins et n'est pas introduite par une étiquette d'appartenance à un complexe.

L'évaluation terminale n'exige pas non plus de choisir une histoire particulière. Pour chaque continuation \(h\), on applique la composition effective de ses mises à jour. Le résultat conserve la paire formée par \(h\) et sa lecture, avec sa multiplicité. Si \(k<\ell<N\), les continuations se décomposent en une union disjointe de préfixes jusqu'à \(\ell\) et de leurs continuations jusqu'à \(N\). La proposition~\ref{viii:rz-num-05-1} démontre qu'évaluer directement ou évaluer ces deux parties produit la même lecture. On obtient ainsi le carré de naturalité du terminal.

Les modules de chemins, les incidences, le transport de mémoire et leurs différentielles procèdent des mêmes prolongements. Leur assemblage ne consiste pas à ajouter à un état cinq coordonnées qui seraient conservées par définition : leurs relations sont construites par des sommes, des bords et des compositions effectives, et sont vérifiées sur leurs générateurs.

\subsection{Déterminants, limite et lecture arithmétique}
Dans chaque troncature finie, le complexe libre admet sa droite déterminante graduée

\[
\operatorname{Det}\Gamma=
\bigotimes_r\left(\bigwedge^{\mathrm{rk}\Gamma_r}\Gamma_r\right)^{(-1)^r}.
\]
La droite est définie par les modules et leurs degrés, aussi bien sur \(\mathbb Z\) qu'après réduction à \(A_N\) ; l'exposant négatif désigne le module dual. Une torsion numérique non nulle exige en outre de spécifier le complexe sur un corps, les bases et les identifications d'homologie pertinentes. Pour calculer les facteurs de Smith, on utilise la matrice entière produite avant la réduction, et non un relèvement arbitraire d'une matrice modulaire. Un bloc carré entier de déterminant non nul est inversible sur \(\mathbb Q\) ; sa réduction est inversible sur \(A_N\) seulement lorsque ce déterminant est une unité, c'est-à-dire lorsqu'il n'est pas divisible par \(3\). Le raffinement des mémoires et des incidences permet de formuler et de vérifier la stabilité des facteurs invariants pertinents ; ces conditions ne se déduisent pas de l'existence de la droite déterminante.

Le passage à la limite conserve les systèmes projectifs compatibles de résidus, les histoires et leurs multiplicités. Les carrés prouvés dans les sections précédentes demeurent commutatifs parce que leurs applications sont données par les mêmes formules à chaque niveau. Pour une coordonnée archimédienne, l'étape supplémentaire consiste à publier des intervalles cylindriques non vides, compatibles par inclusion et de diamètre tendant vers zéro. Leurs fermetures compactes emboîtées ont une intersection réduite à un point : l'existence s'obtient par compacité et l'unicité par la borne du diamètre. C'est la valeur de la lecture limite. Si les intervalles sont semi-ouverts, le point peut n'appartenir qu'à leurs fermetures ; la règle de représentation de frontière détermine alors son écriture positionnelle. La valeur ne remplace pas l'histoire qui l'a produite. L'autre publication conserve les incidences du même état. Les deux destinations partagent leur origine sans pour autant recevoir la même information.

Le secteur des formes normales étudié dans le chapitre arithmétique s’obtient au moyen de la publication normalisée de préfixes du TPK. Son évaluation entière et ses opérations sont des lecteurs ultérieurs sur le substrat construit ici. Les fonctionnelles spectrales, les parties finies et les horloges irréductibles s’appliquent à ces modes publiés. L’identification d’une fonctionnelle concrète à une forme analytique externe requiert à son tour d’écrire son application sur les générateurs et de vérifier leurs produits scalaires : les identités de conservation du présent chapitre ne remplacent pas cette composition.

\clearpage
\subsection{Provenance, vérification et portée}
Ce développement réunit l’architecture de l’auteur des propriétaires 04b1, 04b2 et 04b3 de la structure discrète, ainsi que les identités de bilan, d’incidence et de mémoire des développements 05. Il conserve conjointement les cinq constructions du continu comme opérations sur un même arbre. L’exposé développe leurs preuves locales et explicite les conditions de chaque spécialisation : transport isométrique, poids uniformes, extinction du terminal et non-singularité d’un bloc déterminantal.

Le vérificateur joint confronte des instances exactes des cocycles, de l’annulation pondérée, de la suite rectangulaire et du complexe de mémoire. Ces contrôles accompagnent les démonstrations générales du texte. Leur portée est locale : ni une énumération finie ni la conservation des fichiers de provenance ne sont présentées comme preuve de la positivité de Weil ou comme une identification de l’opérateur spectral qui n’est pas encore composée.
